\ifdefined\gkver\else\def\gkver{student}\fi
\documentclass[\gkver]{gaokaozhenti}
\begin{document}

\chapter{2004年北京春季理综}
%% paperType: 理综物理部分

\begin{enumerate}
\item
%% number: 15
%% paperName: 2004年北京春季理综第 15 题 6 分
%% typeId: 1
%% type: 单选题
%% chapter: 原子和原子核
%% point: 天然放射性
%% method: 无
%% score: 6
%% degree: 750
%% duplicateId: 0
%% body:
钍核 \ce{^{232}_{90}Th} 经过 6 次 $\alpha$ 衰变和 4 次 $\beta$ 衰变后变成铅核，则\xzanswer{C}
\fourchoices[answer=C]
{铅核的符号为 \ce{^{208}_{82}Pb}，它比 \ce{^{232}_{90}Th} 少 8 个中子}
{铅核的符号为 \ce{^{204}_{78}Pb}，它比 \ce{^{232}_{90}Th} 少 16 个中子}
{铅核的符号为 \ce{^{208}_{82}Pb}，它比 \ce{^{232}_{90}Th} 少 16 个中子}
{铅核的符号为 \ce{^{220}_{78}Pb}，它比 \ce{^{232}_{90}Th} 少 12 个中子}
%% answer: C
%% memo:
\memoanswer{
由核反应方程 \ce{^{232}_{90}Th -> 6 ^{4}_{2}He + 4 ^{0}_{-1}e + ^{208}_{82}Pb}，由质量数、电荷数守恒可知铅核的符号为 \ce{^{208}_{82}Pb}，\ce{^{232}_{90}Th} 的中子数为 142 个，\ce{^{208}_{82}Pb} 的中子数为 126 个，比 \ce{^{232}_{90}Th} 少 16 个中子，故 C 正确。
}

\item
%% number: 16
%% paperName: 2004年北京春季理综第 16 题 6 分
%% typeId: 1
%% type: 单选题
%% chapter: 光学
%% point: 光的电磁说
%% method: 无
%% score: 6
%% degree: 750
%% duplicateId: 0
%% body:
对于某单色光，玻璃的折射率比水的大，则此单色光在玻璃中传播时\xzanswer{C}
\fourchoices[answer=C]
{其速度比在水中的大，其波长比在水中的长}
{其速度比在水中的大，其波长比在水中的短}
{其速度比在水中的小，其波长比在水中的短}
{其速度比在水中的小，其波长比在水中的长}
%% answer: C
%% memo:
\memoanswer{
由 $n=\frac{c}{v}=\frac{\lambda_{0}v}{\lambda v}=\frac{\lambda_{0}}{\lambda}$，可知 $n$ 大的 $v$ 小、$\lambda$ 短，故 C 正确。
}

\item
%% number: 17
%% paperName: 2004年北京春季理综第 17 题 6 分
%% typeId: 1
%% type: 单选题
%% chapter: 运动和力
%% point: 连接体
%% method: 无
%% score: 6
%% degree: 550
%% duplicateId: 0
%% body:
图中 a、b 是两个位于固定斜面上的正方形物块，它们的质量相等。$F$ 是沿水平方向作用于 a 上的外力，已知 a、b 的接触面，a、b 与斜面的接触面都是光滑的。正确的说法是\xzanswer{D}
\onepicture[w=6.5cm]{figs/17.png}
\fourchoices[answer=D]
{a、b 一定沿斜面向上运动}
{a 对 b 的作用力沿水平方向}
{a、b 对斜面的正压力相等}
{a 受到的合力沿水平方向的分力等于 b 受到的合力沿水平方向的分力}
%% answer: D
%% memo:
\memoanswer{
利用排除法，$F$ 的大小不同，作用于 a 上，a、b 的运动情况可能不同，故 A 错误；由弹力的方向可判断 a 对 b 的作用力应沿斜面方向，故 B 错误；由受力分析可知，b 对斜面的正压力等于其重力垂直于斜面方向的分力，由于 $F$ 的存在，a 对斜面的正压力大于其重力垂直于斜面方向的分力，故 C 错误；a、b 的质量相等，运动状态相同，即加速度相同，所以 a、b 所受合力大小相同，即沿水平方向的分力大小相等，故 D 正确。
}

\item
%% number: 18
%% paperName: 2004年北京春季理综第 18 题 6 分
%% typeId: 1
%% type: 单选题
%% chapter: 机械波
%% point: 波的图象
%% method: 无
%% score: 6
%% degree: 650
%% duplicateId: 0
%% body:
一简谐横波在 $x$ 轴上传播，波源振动周期 $T=0.1\Us$，在某一时刻的波形如图所示，且此时 a 点向下运动，则\xzanswer{A}
\onepicture[w=8.0cm]{figs/18.png}
\fourchoices[answer=A]
{波速为 $20\Ums$，波向 $x$ 轴正方向传播}
{波速为 $10\Ums$，波向 $x$ 轴负方向传播}
{波速为 $20\Ums$，波向 $x$ 轴负方向传播}
{波速为 $10\Ums$，波向 $x$ 轴正方向传播}
%% answer: A
%% memo:
\memoanswer{
由波形图可知 $\lambda=2\Um$，又 $T=0.1\Us$，则 $v=\frac{\lambda}{T}=20\Ums$，由 a 质点向下运动，可判断波向 $x$ 轴正方向传播，故 A 正确。
}

\item
%% number: 19
%% paperName: 2004年北京春季理综第 19 题 6 分
%% typeId: 1
%% type: 单选题
%% chapter: 电场
%% point: 电场叠加
%% method: 无
%% score: 6
%% degree: 450
%% duplicateId: 0
%% body:
如图，在正六边形的 a、c 两个顶点上各放一带正电的点电荷，电量的大小都是 $q_{1}$，在 b、d 两个顶点上，各放一带负电的点电荷，电量的大小都是 $q_{2}$，$q_{1}>q_{2}$。已知六边形中心 $O$ 点处的场强可用图中的四条有向线段中的一条来表示，它是哪一条？\xzanswer{B}
\onepicture[w=5.5cm]{figs/19.png}
\fourchoices[answer=B]
{$E_{1}$}
{$E_{2}$}
{$E_{3}$}
{$E_{4}$}
%% answer: B
%% memo:
\memoanswer{
正点电荷在 $O$ 点产生的场强方向由该点电荷指向 $O$ 点，负点电荷在 $O$ 点产生的场强方向由 $O$ 点指向该点电荷，先在 $O$ 点分别画出 a、b、c、d 处点电荷在该处的场强，再利用矢量合成，求出合场强方向，与 $E_{2}$ 同向，故 B 正确。
}

\item
%% number: 20
%% paperName: 2004年北京春季理综第 20 题 6 分
%% typeId: 1
%% type: 单选题
%% chapter: 气体性质
%% point: 气体图象
%% method: 图象法
%% score: 6
%% degree: 450
%% duplicateId: 0
%% body:
一定质量的理想气体处于某一平衡状态，此时其压强为 $p_{0}$，有人设计了四种途径，使气体经过每种途经后压强仍为 $p_{0}$。这四种途径是

①先保持体积不变，降低压强，再保持温度不变，压缩体积

②先保持体积不变，使气体升温，再保持温度不变，让体积膨胀

③先保持温度不变，使体积膨胀，再保持体积不变，使气体升温

④先保持温度不变，压缩气体，再保持体积不变，使气体降温

可以断定\xzanswer{D}
\fourchoices[answer=D]
{①、②不可能}
{③、④不可能}
{①、③不可能}
{①、②、③、④都可能}
%% answer: D
%% memo:
\memoanswer{
由 $\frac{pV}{T}=$ 恒量判断：$V$ 不变时 $p\downarrow$、$T\downarrow$，$T$ 不变时 $V\downarrow$、$p\uparrow$，故①正确；$V$ 不变时 $T\uparrow$、$p\uparrow$，$T$ 不变时 $V\uparrow$、$p\downarrow$，故②正确；$T$ 不变时 $V\uparrow$、$p\downarrow$，$V$ 不变时 $T\uparrow$、$p\uparrow$，故③正确；$T$ 不变时 $V\downarrow$、$p\uparrow$，$V$ 不变时 $T\downarrow$、$p\downarrow$，故④正确。综上 D 正确。
}

\item
%% number: 21
%% paperName: 2004年北京春季理综第 21 题 6 分
%% typeId: 1
%% type: 单选题
%% chapter: 电场
%% point: 带电粒子的轨迹类问题
%% method: 无
%% score: 6
%% degree: 450
%% duplicateId: 0
%% body:
如图，$O$ 是一固定的点电荷，另一点电荷 P 从很远处以初速度 $v_{0}$ 射入点电荷 $O$ 的电场，在电场力作用下的运动轨迹是曲线 MN。a、b、c 是以 $O$ 为中心，$R_{a}$、$R_{b}$、$R_{c}$ 为半径画出的三个圆，$R_{c}-R_{b}=R_{b}-R_{a}$。1、2、3、4 为轨迹 MN 与三个圆的一些交点。以 $|W_{12}|$ 表示点电荷 P 由 1 到 2 的过程中电场力的功的大小，$|W_{34}|$ 表示由 3 到 4 的过程中电场力做的功的大小，则\xzanswer{B}
\onepicture[w=4.6cm]{figs/21.png}
\fourchoices[answer=B]
{$|W_{12}|=2|W_{34}|$}
{$|W_{12}|>2|W_{34}|$}
{P、O 两电荷可能同号，也可能异号}
{P 的初速度方向的延长线与 $O$ 之间的距离可能为零}
%% answer: B
%% memo:
\memoanswer{
由轨迹可知，点电荷受力应指向轨迹弯曲的内侧，故 P、Q 电荷间应是引力，故 CD 错误；由于离 Q 越近场强越强，且 $R_{c}-R_{b}=R_{b}-R_{a}$，则 $U_{12}>2U_{34}$，由 $W=qU$ 可判断 $|W_{12}|>2|W_{34}|$，故 B 正确。
}

\item
%% number: 22
%% paperName: 2004年北京春季理综第 22 题 6 分
%% typeId: 1
%% type: 单选题
%% chapter: 磁场
%% point: 带电粒子在磁场中的运动
%% method: 无
%% score: 6
%% degree: 450
%% duplicateId: 0
%% body:
如图，在 $x>0$、$y>0$ 的空间中有恒定的匀强磁场，磁感强度的方向垂直于 $Oxy$ 平面向里，大小为 $B$。现有一质量为 $m$、电量为 $q$ 的带电粒子，在 $x$ 轴上到原点的距离为 $x_{0}$ 的 P 点，以平行于 $y$ 轴的初速度射出此磁场，在磁场作用下沿垂直于 $y$ 轴的方向射出此磁场。不计重力的影响。由这些条件可知\xzanswer{D}
\onepicture[w=4.6cm]{figs/22.png}
\fourchoices[answer=D]
{不能确定粒子通过 $y$ 轴时的位置}
{不能确定粒子速度的大小}
{不能确定粒子在磁场中运动所经历的时间}
{以上三个判断都不对}
%% answer: D
%% memo:
\memoanswer{
由于带电粒子平行 $y$ 轴射入磁场，垂直于 $y$ 轴射出磁场，说明带电粒子在磁场中的运动轨迹是 $\frac{1}{4}$ 圆周，则由 $R=x_{0}=\frac{mv}{qB}=y_{0}$，可确定 AB 错误；由 $t=\frac{T}{4}=\frac{\pi m}{2qB}$，可判断 C 错误，故 D 正确。
}

\item
%% number: 23
%% paperName: 2004年北京春季理综第 23 题 4 分
%% typeId: 4
%% type: 实验题
%% chapter: 直线运动
%% point: 游标卡尺和螺旋测微仪
%% method: 无
%% score: 4
%% degree: 850
%% duplicateId: 0
%% body:
用一主尺最小分度为 $1\Umm$，游标上有 20 个分度的卡尺测量一工件的长度，结果如图所示。可以读出此工件的长度为 \tkanswer[2.5cm]{$10.405\Ucm$}。
\onepicture[w=7.5cm, align=right]{figs/23a.png}
%% answer: 
%% memo:
\memoanswer{
由图知，主尺读数为 $104\Umm$，游标尺的读数为 $1\times0.05\Umm=0.05\Umm$，故工件的长度为 $104\Umm+1\times0.05\Umm=104.05\Umm=10.405\Ucm$。
}

\item
%% number: 23
%% paperName: 2004年北京春季理综第 23 题 8 分
%% typeId: 4
%% type: 实验题
%% chapter: 机械振动
%% point: 单摆测重力加速度
%% method: 无
%% score: 8
%% degree: 650
%% duplicateId: 0
%% body:
在测量重力加速度的实验中，某同学用一根细线和一均匀小球制成单摆。他已经测得此单摆 20 个周期的时间为 $t$，从悬挂点到小球顶端的线长为 $l$，还需要测量的物理量为 \tkanswer[2.5cm]{小球直径 $d$}。将 $g$ 用测得量表示，可得 $g=$ \tkanswer[4cm]{$\frac{1600\pi^{2}\left(l+\frac{d}{2}\right)}{t^{2}}$}。
%% answer: 
%% memo:
\memoanswer{
由题意知，单摆的周期为 $T=\frac{t}{20}$，又单摆的周期为 $T=2\pi\sqrt{\frac{L}{g}}$，联立解得 $g=\frac{1600\pi^{2}\left(l+\frac{d}{2}\right)}{t^{2}}$。
}

\item
%% number: 23
%% paperName: 2004年北京春季理综第 23 题 4 分
%% typeId: 4
%% type: 实验题
%% chapter: 电路
%% point: 测电动势与内阻
%% method: 无
%% score: 4
%% degree: 550
%% duplicateId: 0
%% body:
测量电源的电动势及内阻的实验电路如图 \ref{2004北京春季理综23b} 所示。图 \ref{2004北京春季理综23c} 中给出的器材有：待测的电源（电动势约为 $4\UV$，内阻约为 $2\UO$），电压表（内阻很大，有 $5\UV$、$15\UV$ 两个量程），电流表（内阻不计，有 $0.1\UA$、$1\UA$ 两个量程），滑线变阻器（阻值范围 $0\sim10\UO$），开关。另有导线若干。试按照图 \ref{2004北京春季理综23b} 中的电路在图 \ref{2004北京春季理综23c} 中画出连线，将器材连接成实验电路（要求正确选择电表量程，以保证仪器的安全并使测量有尽可能高的精确度）。

\onepicture[num=1, label=2004北京春季理综23b, w=5.2cm]{figs/23b.png}

\onepicture[num=2, label=2004北京春季理综23c, w=5.5cm]{figs/23c.png}
%% answer: 
\jdanswer{
连线如图所示
}
%% memo:
\memoanswer{
电源电动势大约为 $4\UV$，所以电压表可确定用 $5\UV$ 量程；电路中能接入的最大电阻为 $12\UO$，最小电流为 $I_{\min}=\frac{E}{R_{\max}}=\frac{4}{12}\UA=0.33\UA$，由此确定电流表选用 $1\UA$ 量程，电路连接如图所示。

\onepicture[w=5.5cm]{figs/23_an.png}
}

\item
%% number: 24
%% paperName: 2004年北京春季理综第 24 题 16 分
%% typeId: 6
%% type: 计算题
%% chapter: 曲线运动
%% point: 人造地球卫星
%% method: 无
%% score: 16
%% degree: 550
%% duplicateId: 0
%% body:
神舟五号载入飞船在绕地球飞行的第 5 圈进行变轨，由原来的椭圆轨道变为距地面高度 $h=342\Ukm$ 的圆形轨道。已知地球半径 $R=6.37\times10^{3}\Ukm$，地面处的重力加速度 $g=10\Umsq$。试导出飞船在上述圆轨道上运行的周期 $T$ 的公式（用 $h$、$R$、$g$ 表示），然后计算周期 $T$ 的数值（保留两位有效数字）。
%% answer: 
\jdanswer{
$T=2\pi\sqrt{\frac{(R+h)^{3}}{R^{2}g}}=5.4\times10^{3}\Us$
}
%% memo:
\memoanswer{
设地球质量为 $M$，飞船质量为 $m$，速度为 $v$，圆轨道的半径为 $r$，由万有引力和牛顿第二定律得
\[ G\frac{Mm}{r^{2}}=m\frac{v^{2}}{r}, \]
地面附近有
\[ G\frac{Mm}{R^{2}}=mg, \]
由已知条件知
\[ r=R+h, \]
又
\[ T=\frac{2\pi r}{v}, \]
解以上各式得 $T=2\pi\sqrt{\frac{(R+h)^{3}}{R^{2}g}}$，代入数值解得 $T=5.4\times10^{3}\Us$。
}

\item
%% number: 25
%% paperName: 2004年北京春季理综第 25 题 18 分
%% typeId: 6
%% type: 计算题
%% chapter: 电磁感应
%% point: 电磁感应电路分析
%% method: 无
%% score: 18
%% degree: 450
%% duplicateId: 0
%% body:
如图，直角三角形导线框 abc 固定在匀强磁场中，ab 是一段长为 $l$、电阻为 $R$ 的均匀导线，ac 和 bc 的电阻可不计，ac 长度为 $\frac{l}{2}$。磁场的磁感强度为 $B$，方向垂直纸面向里。现有一段长度为 $\frac{l}{2}$、电阻为 $\frac{R}{2}$ 的均匀导体杆 MN 架在导线框上，开始时紧靠 ac，然后沿 ab 方向以恒定速度 $v$ 向 b 端滑动，滑动中始终与 ac 平行并与导线框保持良好接触。当 MN 滑过的距离为 $\frac{l}{3}$ 时，导线 ac 中的电流是多大？方向如何？
\onepicture[align=right]{\includesvg{figs/25}}
%% answer: 
\jdanswer{
\begin{enumerate}
	\item
	$I_{ac}=\frac{2Blv}{5R}$
	\item
	由 a 流向 c
\end{enumerate}
}
%% memo:
\memoanswer{
当 MN 滑过的距离为 $\frac{l}{3}$ 时，它与 bc 的接触点为 P，如图所示，由几何关系可知 MP 长度为 $\frac{l}{3}$，MP 切割产生的感应电动势为
\[ E=\frac{1}{3}Blv, \]

\onepicture[w=5.5cm]{figs/25_an.jpg}

MP 段的电阻为 $r=\frac{1}{3}R$，MacP 和 MbP 两电路的并联电阻为
\[ r_{\text{并}}=\frac{\frac{1}{3}\times\frac{2}{3}}{\frac{1}{3}+\frac{2}{3}}R=\frac{2}{9}R, \]
由欧姆定律得，PM 中的电流为
\[ I=\frac{E}{r+r_{\text{并}}}, \]
ac 中的电流为
\[ I_{ac}=\frac{2}{3}I, \]
解得
\[ I_{ac}=\frac{2Blv}{5R}, \]
由右手定则知，MP 中的感应电流的方向由 P 流向 M，所以电流 $I_{ac}$ 的方向由 a 流向 c。
}

\item
%% number: 34
%% paperName: 2004年北京春季理综第 34 题 22 分
%% typeId: 6
%% type: 计算题
%% chapter: 运动和力
%% point: 动量和能量
%% method: 无
%% score: 22
%% degree: 350
%% duplicateId: 0
%% body:
如图，abc 是光滑的轨道，其中 ab 是水平的，bc 为与 ab 相切的位于竖直平面内的半圆，半径 $R=0.30\Um$。质量 $m=0.20\Ukg$ 的小球 A 静止在轨道上，另一质量 $M=0.60\Ukg$、速度 $v_{0}=5.5\Ums$ 的小球 B 与小球 A 正碰。已知相碰后小球 A 经过半圆的最高点 c 落到轨道上距 b 点为 $l=4\sqrt{2}R$ 处，重力加速度 $g=10\Umsq$，求：
\begin{enumerate}
	\item
	碰撞结束后，小球 A 和 B 的速度的大小。
	\item
	试论证小球 B 是否能沿着半圆轨道到达 c 点。
\end{enumerate}
\onepicture[w=8.0cm, align=right]{figs/34.png}
%% answer: 
\jdanswer{
\begin{enumerate}
	\item
	$v_{1}=6\Ums$，$v_{2}=3.5\Ums$
	\item
	由 $v_{2}=3.5\Ums$、$v_{b}=3.9\Ums$ 可知 $v_{2}<v_{b}$，所以小球 B 不能达到半圆轨道的最高点。
\end{enumerate}
}
%% memo:
\memoanswer{
（1）以 $v_{1}$ 表示小球 A 碰后的速度，$v_{2}$ 表示小球 B 碰后的速度，$v_{1}'$ 表示小球 A 在半圆最高点的速度，$t$ 表示小球 A 从离开半圆最高点到落在轨道上经过的时间。由平抛规律得
\[ 4\sqrt{2}R=v_{1}'t, \]①
\[ 2R=\frac{1}{2}gt^{2}, \]②
由机械能守恒得
\[ mg(2R)+\frac{1}{2}mv_{1}'^{2}=\frac{1}{2}mv_{1}^{2}, \]③
由动量守恒得
\[ Mv_{0}=mv_{1}+Mv_{2}, \]④
由①②③④求得
\[ v_{1}=2\sqrt{3Rg}, \qquad v_{2}=v_{0}-2\frac{m}{M}\sqrt{3Rg}, \]
代入数值得 $v_{1}=6\Ums$，$v_{2}=3.5\Ums$。

（2）假设 B 球刚能沿着半圆轨道上升到 c 点，则在 c 点时，轨道对它的作用力等于零。以 $v_{c}$ 表示它在 c 点的速度，$v_{b}$ 表示它在 b 点相应的速度，由牛顿第二定律和机械能守恒得
\[ Mg=M\frac{v_{c}^{2}}{R}, \]
\[ \frac{1}{2}Mv_{c}^{2}+Mg(2R)=\frac{1}{2}Mv_{b}^{2}, \]
联立解得 $v_{b}=\sqrt{5Rg}$，代入数值得 $v_{b}=3.9\Ums$。由 $v_{2}=3.5\Ums$ 可知 $v_{2}<v_{b}$，所以小球 B 不能达到半圆轨道的最高点。
}

\end{enumerate}

\end{document}
